% =========================================================
\chapter{Estado da Arte}
% =========================================================

O Estado da Arte (EA) deve identificar o conhecimento acumulado, as principais abordagens existentes, suas limitações e, sobretudo, a lacuna científica ou tecnológica na qual o trabalho está inserido.
Noutras palavras\textcolor{blue}{, o EA posiciona a solução de um TCC no espaço-tempo das soluções atuais, visando apresentar o que há de mais moderno e comparando-o com a solução ora proposta. }

Para fins de orientaçoes gerais na produção do TCC, deve-se trabalhar conforme as orientações disponibilizadas no sítio do IFCE:
\begin{enumerate}
    \item \textbf{Orientações Gerais para Entrega de Trabalho Final} \\ \url{https://acervo.ifce.edu.br/proen/bibliotecas/entrega-de-trabalhos-academicos}
    
    \item \textbf{PPC do curso de Engenharia de Controle e Automação}\\ \url{https://ifce.edu.br/maracanau/menu/cursos/superiores/bacharelados/controle_automacao/pdf/ppc-eca-1.pdf}
    
\end{enumerate}


% ---------------------------------------------------------
\section{Considerações iniciais}
% ---------------------------------------------------------

\begin{enumerate}
	\item Qual é o domínio científico específico investigado?
	\item Quais são os principais subtemas envolvidos?
	\item Quais critérios foram utilizados para selecionar os trabalhos analisados?
	\item Qual período da literatura foi considerado?
	\item Quais bases de dados foram utilizadas?
	\item Quais palavras-chave foram utilizadas nas buscas?
	\item Quais critérios de inclusão e exclusão foram adotados?
	\item A revisão é narrativa, sistemática, bibliométrica ou uma combinação dessas abordagens?
	\item Quais são as principais \textbf{tendências} identificadas?
	\item Como a revisão conduzirá à identificação da lacuna científica?
\end{enumerate}

% ---------------------------------------------------------
\section{Fronteira do conhecimento nos assuntos envolvidos}
% ---------------------------------------------------------

\begin{enumerate}
	\item Quais são os trabalhos mais relevantes e recentes sobre o problema?
	\item Quais pesquisadores ou grupos são referências no tema?
	\item Quais métodos são predominantes na literatura?
	\item Quais tecnologias ou arquiteturas são atualmente utilizadas?
	\item Quais modelos matemáticos aparecem com maior frequência?
	\item Quais métricas são utilizadas para avaliar as soluções?
	\item Quais resultados são considerados estado da arte?
	\item Quais limitações os próprios autores apontam?
	\item Quais resultados da literatura são conflitantes ou inconclusivos?
	\item Qual lacuna permanece aberta após a análise crítica da literatura?
\end{enumerate}

{Pergunta central.}
\textbf{Qual é o estado atual do conhecimento e exatamente onde meu trabalho se
	posiciona?}

% ---------------------------------------------------------
\section{Visão holística do estado da arte}
% ---------------------------------------------------------

\begin{enumerate}
	\item Como os diferentes trabalhos da literatura se relacionam?
	\item Quais são as principais linhas de pesquisa?
	\item Quais abordagens podem ser agrupadas por método?
	\item Quais abordagens podem ser agrupadas por aplicação?
	\item Quais variáveis ou métricas permitem comparar os trabalhos?
	\item Quais são as vantagens e desvantagens de cada abordagem?
	\item Quais soluções são complementares?
	\item Quais soluções são concorrentes?
	\item Como os trabalhos evoluíram historicamente?
	\item Qual espaço científico permanece disponível para a solução proposta?
\end{enumerate}

Sugestão de síntese comparativa.

\begin{center}
	\begin{tabular}{p{2.5cm}p{2cm}p{2cm}p{2cm}p{2cm}p{3cm}}
		\toprule
		\textbf{Trabalho} &
		\textbf{Método} &
		\textbf{Aplicação} &
		\textbf{Dados} &
		\textbf{Métrica} &
		\textbf{Principal resultado/limitação} \\
		\midrule
		Autor 1 & Método & Aplicação & Dados & Métrica & Resultado e limitação \\
		Autor 2 & Método & Aplicação & Dados & Métrica & Resultado e limitação \\
		Autor 3 & Método & Aplicação & Dados & Métrica & Resultado e limitação \\
		\bottomrule
	\end{tabular}
\end{center}

\textcolor{blue}{Note que isto traz destaque amplo para o que será proposto.}
% ---------------------------------------------------------
\section{Considerações finais}
% ---------------------------------------------------------

\begin{enumerate}
	\item Quais foram as principais conclusões do Estado da Arte?
	\item Quais abordagens apresentaram melhor desempenho?
	\item Quais limitações foram identificadas?
	\item Quais lacunas permanecem?
	\item Qual lacuna será atacada pelo presente trabalho?
	\item Por que essa lacuna é relevante?
	\item Como a solução proposta se posiciona em relação à literatura?
	\item Qual hipótese emerge da revisão?
	\item Quais requisitos a solução proposta deve atender?
	\item Como o Estado da Arte justifica o capítulo seguinte?
\end{enumerate}

\textcolor{blue}{O redator do TCC deve atentar a essas perguntas, pois a banca de avaliação haverá de realizá-las sobre o trabalho por ocasião da defesa.}
